\scanpage{96}
% archive foliation: 95
\hfill n\textsuperscript{o} 285
\begin{center}
RAPPORT sur la RÉÉDITION\\
de\\
TOPOLOGIE GENERALE\\
Chap.III
\end{center}

§2 ; les n\textsuperscript{os} 4 et 8 sont remaniés de la façon suivante :

4. - \uline{Espaces homogènes}

Soient $E$ un espace topologique, $G$ un groupe topologique, et supposons $E$ muni du groupe d'opérateurs $G$ (Alg. chap.I, §7, n\textsuperscript{o}2) ; nous dirons que $G$ \uline{opère continûment} dans $E$ si l'application $(s,x)\longrightarrow s.x$ de $G\times E$ dans $E$ est \uline{continue}. Lorsqu'il en est ainsi, pour tout $s\in G$, l'application $x\longrightarrow s.x$ de $E$ dans $E$ est un \uline{homéomorphisme} de $E$ sur lui-même, car c'est une bijection continue dont la bijection réciproque $x\longrightarrow s^{-1}.x$ est aussi continue.

Supposons en outre que $G$ opère \uline{transitivement} dans $E$, autrement dit que $E$ soit (au sens algébrique) un \uline{espace homogène} de $G$ (Alg., chap.I, §7,n\textsuperscript{o}6). Soient $x_{\unsure{o|0}}$ un point de $E$, $H_{\unsure{o|0}}$ le sous-groupe de $G$ formé des opérateurs laissant invariant $x_{\unsure{o|0}}$ ( en \uline{Alg.}, Bourbaki n'a pas cru bon de donner un nom à ce sous-groupe ; faut-il rattraper cela, et l'appeler ``stabilisateur'' de $x_{\unsure{o|0}}$ ?). La surjection continue $s\longrightarrow s.x_{\unsure{o|0}}$ de $G$ sur $E$ se factorise canoniquement en
\diagram{\begin{tikzcd}[ampersand replacement=\&]
G\arrow[r,"f"]\&G/H_{\unsure{o|0}}\arrow[r,"g"]\&E
\end{tikzcd}}
où $f$ est l'application canonique de $G$ sur l'espace homogène $G/H_{\unsure{o|0}}$ et $g$ la bijection $s.H_{\unsure{o|0}}\longrightarrow s.x_{\unsure{o|0}}$ de $G/H_{\unsure{o|0}}$ sur $E$ ; on sait en outre (Chap.I, 2\textsuperscript{eme} éd. §9, n\textsuperscript{o}2, th.1) que si on munit $G/H_{\unsure{o|0}}$ de la topologie quotient de celle de $G$ par la relation d'équivalence $s^{-1}t\in H_{\unsure{o|0}}$, $g$ est \uline{continue}. On dit que $E$ est un \uline{espace homogène topologique}(correspondant
